\documentclass[../../main2.tex]{subfiles}

\begin{document}

\chapter{多因分账：当原因搅在一起}
\label{ch:multiple}

\begin{motivation}
	上一章遗留的问题：无你检验在投毒案里出了洋相——甲乙各投致死量，拿掉谁都还死，难道都无责？单人剧本装不下真实世界：事故、革命、爆款，从来是多因同台。\\
	\textbf{核心动机}：把「各占多少」从单因推进到多因——交互项必须显式处理，分账方案必须挑明代价。这是全书最硬的一块，也是「分锅」这门手艺的心脏。
\end{motivation}

\begin{logicmap}
\begin{tikzpicture}[node distance=0.85cm, every node/.style={draw, rectangle, rounded corners, align=center, font=\small}]
\node (q) {拿掉谁都还发生，谁的锅？};
\node (c) [below=of q] {双充分与合力：两种难案};
\node (d) [below=of c] {交互项：一加一不等于二};
\node (p) [below=of d] {两套分账：平摊 vs 按边际};
\node (n) [below=of p] {份额与量级：归一化陷阱};
\draw[-{Stealth}] (q) -- (c);
\draw[-{Stealth}] (c) -- (d);
\draw[-{Stealth}] (d) -- (p);
\draw[-{Stealth}] (p) -- (n);
\end{tikzpicture}
\end{logicmap}

\section{两种难案：拿掉谁都还发生}

\begin{readingnote}
	你有没有想过，两个人同时放火、同一栋楼烧了，法律凭什么都判？\\
	按上一章的无你检验，两把火都「拿掉也烧」——漏洞出在哪？
\end{readingnote}

两把火的经典案型（法学院一年级就吵）：甲、乙各自独立放火，火势都足以烧光全楼，楼真的烧光了。无你检验说：拿掉甲，乙也会烧掉——甲无责；拿掉乙，甲也会烧掉——乙无责。结论「都无责」荒唐，说明\textbf{检验本身}在这里不够用。

另一种难案反过来：甲的毒单独不致死，乙的毒单独也不致死，\textbf{合起来}致死。拿掉谁都无害——但谁都脱不了干系。这叫合力原因。

\begin{tcolorbox}[colback=green!5, title=\textit{生活类比}]
	双充分像两把都能开同一扇门的钥匙，合力像两把各开半扇的钥匙。前者是「抢功」，后者是「拼单」。生活里爆款常是后者：好剧本 + 好档期 + 好宣发，缺一不可。
\end{tcolorbox}

当下的人群里，这样的案子不少：

\begin{itemize}
	\item \textbf{上班族}：项目黄了——预算砍半（甲）和核心离职（乙），任何一件单独都够呛，合起来致命。复盘会上互指「没有你我也黄」。
	\item \textbf{学生}：考研失利——那场病（甲）加上缩招（乙），单独都可能扛过，一起来就崩了。
	\item \textbf{历史}：崇祯的明朝——天灾、流民、边患、财政，哪条是「那个原因」？单拿掉一条，亡国还会发生吗？
\end{itemize}

\paragraph*{逻辑衔接}
	两类难案认出来了。但「合力」和「抢功」的账还是一笔糊涂账。下一节把糊涂账里最肥的那块——交互项——拽到明处。

\section{交互项：一加一不等于二}
\label{sec:multiple-interaction}

先看案例（铁律：公式前站案例）。

\begin{examplebox}
	\textbf{案例：工地事故}。事故调查：防护网缺失（A），加上违规堆载（B）。统计三个时期——\\
	都有：事故率 $0.30$；只有 A 没 B：$0.12$；只有 B 没 A：$0.08$；都没有：$0.02$。\\
	问：A、B 各占多少？
\end{examplebox}

照第 8 章的算法各算各的：$C_A = 0.30 - 0.08 = 0.22$，$C_B = 0.30 - 0.12 = 0.18$。加起来 $0.40$。但总抬升是 $0.30 - 0.02 = 0.28$。\\
凭空多算了 $0.12$——这 $0.12$ 不是谁的功劳，是 A、B 功能重叠、被重复记账的部分——交互项（interaction）。

\begin{mathbox}{交互项：多出来（或少掉）的那块}
	\[
	I_{AB} \;=\; \big[P(D \mid A,B) - P(D \mid \neg A, \neg B)\big] - C_A - C_B
	\]
	\textbf{代入案例}：$I_{AB} = 0.28 - 0.22 - 0.18 = -0.12$。\\
	负的！为什么？因为防护网和堆载\textbf{功能重叠}：挡住一个，另一个危害就小——互相「抢锅」。反过来，毒药案的 $I_{AB}$ 是正的：协同增效。\\[0.5em]
	\textbf{为什么这样定义好？}它强迫你面对一个事实：分账不是切蛋糕，蛋糕是蛋糕，\textbf{账本上可能凭空多出或少掉一块}。不显式处理 $I_{AB}$ 的分账，都是和稀泥。
\end{mathbox}

\mathlink{app:interaction}{完整推导与符号表见数学附录 A.4。}

\begin{warningbox}
	\textbf{交互项不许随便分掉。}$I_{AB}$ 是「合谋」的部分，既不姓 A 也不姓 B。偷偷塞给谁，谁就背黑锅。怎么处置它，正是下一节两套方案的分歧点。
\end{warningbox}

\paragraph*{逻辑衔接}
	交互项明了。现在剩下的就是分账方案：这块「合谋」的账，摊给谁？两套主流方案，代价各不相同。

\section{两套分账：平摊，还是按边际}

\begin{readingnote}
	你有没有想过，四个人凑钱创业、公司做成了，股份凭什么这么分？\\
	是按人头平摊，还是按「没有你行不行」？两种方案都有人拥护——因为\textbf{它们答的不是同一个问题}。
\end{readingnote}

\textbf{方案一：夏普利值}（Shapley value）——对称平摊。思路：把所有「上场顺序」都排一遍，算每个原因在所有顺序里的平均边际贡献，交互项按人次均摊。优点：公平感强，谁也不吃亏；缺点：把「合谋」稀释了，会掩盖主犯。

\textbf{方案二：边际贡献排序}——只问「拿掉你会怎样」。优点：主犯立刻现形，和无你检验无缝衔接；缺点：加不起来——各算各的，份额合计可能超过总量，无法当「百分比」用。

怎么选？我给一条可操作的规矩：

\begin{itemize}
	\item 要\textbf{分一笔确定的钱、必须加总}（赔偿、股份、奖金）——用夏普利值，账面能合拢；
	\item 要\textbf{抓主因、定优先级}（事故整改、政策排序）——用边际贡献，别让稀释掩盖杠杆点。
\end{itemize}

\begin{examplebox}
	\textbf{案例：分摊四百万赔偿}。事故有 A、B、C 三个原因，交互部分 60 万。\\
	夏普利：交互的 60 万按三方均摊，各背 20 万，加上各自主贡献——四百万分文不差。\\
	边际：A 是「拿掉就不出事」的主因，先给 A 记满 220 万；B、C 的份额相加只剩 120 万，再分交互的 60 万就冒了。\\
	结论：赔偿用前者，整改清单用后者。\end{examplebox}

\lowfreq{诚实标注：两套方案都是约定，不是定理。哪种「对」，取决于你要回答什么问题——这本身是不可通约性（第 2 章）在因果空间的回声。}

\begin{questionbox}
	\textit{现在你可能想反驳：「按你这么说，责任数字是人定的？那还有什么客观可言？」}\\
	\textbf{分账规则是约定，但约束是客观的。}就像「米」的定义是约定，但楼塌不塌不看你用不用米。方案可挑，数据、因果方向、交互项不容挑——选了方案就得承担它的代价，不许装无辜。
\end{questionbox}

\paragraph*{逻辑衔接}
	账分完了，数字出来了。但数字的读法有个大坑——下一节，也是本章最后一个坑：把账读成百分比之后，你会丢掉什么。

\section{归一化的陷阱：份额与量级}

\begin{readingnote}
	你有没有想过，为什么有的事故报告说「A 占四成」，你听完觉得事不大；有的说「A 占四成」，你听完睡不着？\\
	同一个百分比，背后的\textbf{总量}可能差一百倍。
\end{readingnote}

把贡献除以贡献之和，得到\textbf{份额表}——好读，能进新闻标题。但归一化（normalization）扔掉了绝对量级：

\begin{center}
\begin{tabular}{lccc}
\toprule
场景 & $C_A$ & $C_B$ & 份额（A : B） \\
\midrule
小事故 & 0.04 & 0.06 & 四成 : 六成 \\
大灾难 & 0.40 & 0.60 & 四成 : 六成 \\
\bottomrule
\end{tabular}
\end{center}

两行的份额表\textbf{一模一样}，风险差十倍。只报份额不报量级，等于只报「他占四成」不说「总共多大」。通胀里也一样：物价「结构性上涨」的份额图永远好看，涨幅的绝对数字才是钱包。

\begin{warningbox}
	\textbf{读贡献表的规矩：份额和量级必须成对出现。}凡是只给你百分比、不给总量的分析——营销话术、甩锅报告、大势判断——先问一句：分母呢？
\end{warningbox}

\begin{reader_anticipation}
	\item 交互项既然「既不姓 A 也不姓 B」，现实中法官到底怎么判的？→ 判决书用「主次责」粗粒度处理，本书把粗粒度还原成数——附录 A.4 有完整账。
	\item 账上的 A、B、C 是字母，社会里的「因素」长什么样？→ Part IV 把字母换成齿轮：观念、经济、权力、法律。
\end{reader_anticipation}

\paragraph*{逻辑衔接}
	到此，分账的手艺齐了：无你检验给单因，交互项管合谋，两套方案分场景，份额量级成对读。但账上的 A、B、C 还是抽象字母——社会里的「因素」到底长什么样？Part IV 把字母换成齿轮：观念、经济、权力、法律。

\begin{thinkbox}
\begin{enumerate}
	\item 找一次你经历过的「多人背锅」场景（小组作业、项目失败、家务分工）：试着用「边际贡献」排序每个人，再用「夏普利式平摊」分一次——两次结果差在哪？哪个更接近当时实际的处理？
	\item \textbf{反事实}：投毒案里，如果乙的毒比甲晚下十分钟，两人的账会变吗？「上场顺序」变了，夏普利式的平均贡献会怎么动？
	\item \textbf{反事实}：如果防护网和违规堆载的功能\textbf{完全重叠}（挡一个等于挡两个），交互项会变成多少？谁该背主锅？——用本章的公式推一遍，再说说你的直觉和公式打不打架。
\end{enumerate}
\end{thinkbox}

\end{document}
